\documentclass[11pt]{article}
\usepackage[T1]{fontenc}
\usepackage{textcomp}
\usepackage[margin=0.75in]{geometry}
\usepackage{amsmath,amssymb,amsthm}
\usepackage{array,booktabs}
\usepackage{hyperref}
\setlength{\parindent}{0pt}
\newtheorem{theorem}{Theorem}
\theoremstyle{definition}
\newtheorem{definition}{Definition}
\title{Flow Proof Artifact: Monoid}
\author{Generated by \texttt{flow doc proof}}
\date{Flow Proof Book}
\begin{document}
\maketitle
Monoid axioms for an associative operation with identity.\\[0.5em]
\textbf{Source.} dummit-foote — *Abstract Algebra*, §1.1\\[0.5em]
\bigskip
\noindent{\large \textbf{Definition 1.} \textit{Monoid multiplication associates} \hfill Monoid \cdot multiplication \cdot parentheses do not matter}\par
\smallskip\noindent\textit{Coordinate: Monoid \cdot multiplication \cdot parentheses do not matter}\par
\smallskip\noindent\textit{Source: dummit-foote}\par
\medskip
\noindent\textbf{Goal.} Monoid multiplication associates\par
\noindent$\displaystyle \forall a \in Monoid \forall b \in Monoid \forall c \in Monoid\quad (a \cdot b) * c = a * (b \cdot c)$\par
\medskip
\noindent
\begin{tabular}{@{} >{\raggedright\arraybackslash}p{0.06\textwidth} >{\raggedright\arraybackslash}p{0.47\textwidth} >{\raggedleft\arraybackslash}p{0.41\textwidth} @{}}
\textbf{\#} & \textbf{Proof} & \textbf{Mathematics} \\
\midrule
\textbf{1.} & We stipulate parentheses do not matter for multiplication on Monoid — this is a definition, not a derived fact. & \\[0.45em]
\textbf{2.} & This follows directly from the definition: (a times b) times c equals a times (b times c). Hence proven. & $\displaystyle (a \cdot b) * c = a * (b \cdot c)$ \\[0.45em]
\bottomrule
\end{tabular}
\label{thm:Monoid-multiplication-parentheses_do_not_matter}
\par\medskip\hrule
\bigskip
\noindent{\large \textbf{Definition 2.} \textit{The identity is a left unit} \hfill Monoid \cdot identity \cdot one is the left identity}\par
\smallskip\noindent\textit{Coordinate: Monoid \cdot identity \cdot one is the left identity}\par
\smallskip\noindent\textit{Source: dummit-foote}\par
\medskip
\noindent\textbf{Goal.} The identity is a left unit\par
\noindent$\displaystyle \forall m \in Monoid\quad 1 \cdot m = m$\par
\medskip
\noindent
\begin{tabular}{@{} >{\raggedright\arraybackslash}p{0.06\textwidth} >{\raggedright\arraybackslash}p{0.47\textwidth} >{\raggedleft\arraybackslash}p{0.41\textwidth} @{}}
\textbf{\#} & \textbf{Proof} & \textbf{Mathematics} \\
\midrule
\textbf{1.} & We stipulate one is the left identity for identity on Monoid — this is a definition, not a derived fact. & \\[0.45em]
\textbf{2.} & This follows directly from the definition: 1 times m equals m. Hence proven. & $\displaystyle 1 \cdot m = m$ \\[0.45em]
\bottomrule
\end{tabular}
\label{thm:Monoid-identity-one_is_the_left_identity}
\par\medskip\hrule
\bigskip
\noindent{\large \textbf{Definition 3.} \textit{The identity is a right unit} \hfill Monoid \cdot identity \cdot one is the right identity}\par
\smallskip\noindent\textit{Coordinate: Monoid \cdot identity \cdot one is the right identity}\par
\smallskip\noindent\textit{Source: dummit-foote}\par
\medskip
\noindent\textbf{Goal.} The identity is a right unit\par
\noindent$\displaystyle \forall m \in Monoid\quad m \cdot 1 = m$\par
\medskip
\noindent
\begin{tabular}{@{} >{\raggedright\arraybackslash}p{0.06\textwidth} >{\raggedright\arraybackslash}p{0.47\textwidth} >{\raggedleft\arraybackslash}p{0.41\textwidth} @{}}
\textbf{\#} & \textbf{Proof} & \textbf{Mathematics} \\
\midrule
\textbf{1.} & We stipulate one is the right identity for identity on Monoid — this is a definition, not a derived fact. & \\[0.45em]
\textbf{2.} & This follows directly from the definition: m times 1 equals m. Hence proven. & $\displaystyle m \cdot 1 = m$ \\[0.45em]
\bottomrule
\end{tabular}
\label{thm:Monoid-identity-one_is_the_right_identity}
\par\medskip\hrule
\end{document}
